\documentclass[11pt,a4paper]{article}
\usepackage[utf8]{inputenc}
\usepackage[T1]{fontenc}
\usepackage{lmodern}
\usepackage[french]{babel}
\usepackage{amsmath,amssymb,mathtools}
\usepackage{icomma}
\usepackage{xcolor}
\usepackage{colortbl}
\usepackage{tikz}
\usepackage[most]{tcolorbox}
\usepackage{enumitem}
\usepackage{array}
\usepackage{booktabs}
\usepackage{fancyhdr}
\usepackage[hidelinks]{hyperref}
\usetikzlibrary{arrows.meta,positioning,calc,patterns,backgrounds,shapes.geometric,decorations.pathreplacing,decorations.markings}
\usepackage[a4paper,margin=2.15cm,headheight=26pt,headsep=12pt]{geometry}
\usepackage{titlesec}

% ---------------- Palette ----------------
\definecolor{primary}{RGB}{0,151,169}
\definecolor{primarydark}{RGB}{0,88,101}
\definecolor{defcol}{RGB}{0,114,178}
\definecolor{thmcol}{RGB}{0,140,120}
\definecolor{formulacol}{RGB}{198,93,9}
\definecolor{excol}{RGB}{0,135,90}
\definecolor{attncol}{RGB}{200,40,55}
\definecolor{methcol}{RGB}{90,90,100}
\definecolor{remarkcol}{RGB}{120,94,200}
\definecolor{barcol}{RGB}{0,151,169}
\definecolor{barcold}{RGB}{0,110,124}
\definecolor{modecol}{RGB}{200,55,55}
\definecolor{meancol}{RGB}{0,90,200}
\definecolor{medcol}{RGB}{160,90,190}
\definecolor{quartcol}{RGB}{0,135,90}
\definecolor{ogivecol}{RGB}{176,74,0}

% ---------------- Boîtes ----------------
\tcbset{boxbase/.style={enhanced,breakable,boxrule=0.9pt,arc=2.5pt,
  left=3mm,right=3mm,top=2.3mm,bottom=2.3mm,
  fonttitle=\bfseries,coltitle=white,toptitle=0.6mm,bottomtitle=0.6mm}}

\newtcolorbox{definition}[1][]{boxbase,colback=defcol!6,colframe=defcol,
  title={\textsc{Définition}\ifx\relax#1\relax\relax\else\textmd{ : \itshape #1}\fi}}
\newtcolorbox{propriete}[1][]{boxbase,colback=thmcol!7,colframe=thmcol,
  title={\textsc{Propriété}\ifx\relax#1\relax\relax\else\textmd{ : \itshape #1}\fi}}
\newtcolorbox{formule}[1][]{boxbase,colback=formulacol!7,colframe=formulacol,
  title={\textsc{Formule}\ifx\relax#1\relax\relax\else\textmd{ : \itshape #1}\fi}}
\newtcolorbox{exemple}[1][]{boxbase,colback=excol!6,colframe=excol,
  title={\textsc{Exemple}\ifx\relax#1\relax\relax\else\textmd{ : \itshape #1}\fi}}
\newtcolorbox{methode}[1][]{boxbase,colback=methcol!8,colframe=methcol,sharp corners,
  title={\textsc{Méthode}\ifx\relax#1\relax\relax\else\textmd{ : \itshape #1}\fi}}
\newtcolorbox{remarque}[1][]{boxbase,colback=remarkcol!7,colframe=remarkcol,
  title={\textsc{Astuce}\ifx\relax#1\relax\relax\else\textmd{ : \itshape #1}\fi}}
\newtcolorbox{attention}[1][]{boxbase,colback=attncol!7,colframe=attncol,
  title={\textsc{Attention}\ifx\relax#1\relax\relax\else\textmd{ : \itshape #1}\fi}}

% Boîte "exercice" et "corrigé" (titre libre passé en argument)
\newtcolorbox{exobox}[1]{boxbase,colback=primary!4,colframe=primary,
  before skip=8pt,after skip=8pt,title={#1}}
\newtcolorbox{corbox}[1]{boxbase,colback=excol!5,colframe=excol!85!black,
  before skip=8pt,after skip=8pt,title={#1}}

% ---------------- En-tête ----------------
\newcommand{\docpart}[1]{\def\thedocpart{#1}}
\docpart{}
\pagestyle{fancy}
\fancyhf{}
\fancyhead[L]{\small\textcolor{primarydark}{\textbf{Fiche 8} — Statistiques descriptives}}
\fancyhead[R]{\small\textcolor{primarydark}{\textsc{\thedocpart}}}
\fancyfoot[C]{\small\thepage}
\renewcommand{\headrulewidth}{0.8pt}
\renewcommand{\headrule}{\hbox to\headwidth{\color{primary}\leaders\hrule height \headrulewidth\hfill}}

\titleformat{\section}{\large\bfseries\color{primarydark}}{\thesection}{0.6em}{}
\titleformat{\subsection}{\normalsize\bfseries\color{primary}}{\thesubsection}{0.6em}{}
\setlist{topsep=2pt,itemsep=1.5pt,parsep=0pt}

% Bandeau de titre du document
\newcommand{\bandeau}[2]{%
\begin{tcolorbox}[boxbase,colback=primary,colframe=primarydark,halign=center,
   before skip=2pt,after skip=12pt]
{\LARGE\bfseries\color{white} #1}\\[3pt]{\normalsize\color{white} #2}
\end{tcolorbox}}

\newcommand{\ecm}{\sigma_m}
\newcommand{\ect}{\sigma}
\newcommand{\med}{M\!e}
\docpart{Moyen — Thème 2}
\begin{document}
\bandeau{Thème 2 — Mode et moyenne}{Niveau Moyen — exercices}

\noindent\textit{Objectif : mener un calcul de moyenne complet et exploiter la linéarité dans les
deux sens.}\medskip

\begin{exobox}{Exercice 1 — Moyenne d'une série en classes}
Masse (g) de $N=40$ souris :
\begin{center}
\begin{tabular}{|c|c|c|c|c|c|}
\hline
Classes (g) & $[20,24[$ & $[24,28[$ & $[28,32[$ & $[32,36[$ & $[36,40[$ \\\hline
$n_i$ & 4 & 10 & 14 & 8 & 4 \\\hline
\end{tabular}
\end{center}
Calculer la masse moyenne (utiliser les centres de classe).
\end{exobox}

\begin{exobox}{Exercice 2 — Moyenne d'une série discrète}
Nombre d'enfants par famille dans $25$ familles :
\begin{center}
\begin{tabular}{|c|c|c|c|c|c|}
\hline
$x_i$ & 0 & 1 & 2 & 3 & 4 \\\hline
$n_i$ & 4 & 9 & 8 & 3 & 1 \\\hline
\end{tabular}
\end{center}
Calculer le nombre moyen d'enfants par famille.
\end{exobox}

\begin{exobox}{Exercice 3 — Changement d'unité (linéarité)}
Une série de températures a pour moyenne $m=20\ (^\circ\mathrm{C})$. On les convertit en degrés
Fahrenheit par la formule $F=1{,}8\,C+32$. Que vaut la moyenne en $^\circ$F ?
\end{exobox}

\begin{exobox}{Exercice 4 — Retrouver la transformation}
La masse moyenne d'un lot est $m=50$\,kg. On souhaite obtenir une moyenne de $55$\,kg.
\begin{enumerate}[label=\textbf{(\alph*)}]
  \item Quelle constante faut-il ajouter à chaque masse ?
  \item Par quel facteur $\alpha$ faudrait-il plutôt \emph{multiplier} chaque masse (sans rien ajouter) ?
\end{enumerate}
\end{exobox}

\begin{exobox}{Exercice 5 — Mode et moyenne ensemble}
\begin{center}
\begin{tabular}{|c|c|c|c|c|}
\hline
Classes & $[10,20[$ & $[20,30[$ & $[30,40[$ & $[40,50[$ \\\hline
$n_i$ & 6 & 15 & 12 & 7 \\\hline
\end{tabular}
\end{center}
Déterminer la classe modale et la moyenne. Le mode et la moyenne coïncident-ils ?
\end{exobox}
\end{document}
